\section{Conclusion}
Our experiments show that NetTrack demonstrates high universality across diverse outdoor scenes, whereas the OpenCV‑based Birds’ Track performs satisfactorily only in simple scenarios.

In terms of efficiency, the OpenCV method runs quickly on CPU alone, while NetTrack requires longer runtime and higher memory even on CPU, and depends on substantial GPU resources for training.

Therefore, for simple scenes with limited computation and memory, we recommend Birds’ Track.

These work has several limitations: (1) experiments were confined to the VAL set due to time and memory constraints; (2) detecting birds against dynamic backgrounds (e.g., water surfaces) remains challenging; (3) the soft‑weight parameter for specular masks, though balanced in tested videos, may need per‑video adjustment for unseen datasets; (4) comprehensive ablation studies and comparative curves are lacking.

Future work will involve larger datasets, detailed ablation analyses, and further optimization of the Birds’ Track framework to enhance its efficiency and broader applicability.